\section{Conclusion}

This study shows that accelerating nonlinear GPU SVM training requires more
than a faster RBF kernel. Kernel-column reuse, fused dataflow, and a coherent
FP16 operator reduce redundant computation and memory traffic, producing a
1.351$\times$ matched speedup on EMBER-65K. A separate held-out evaluation
detects no statistically significant accuracy difference in the measured
experiment. The numerical solution nevertheless differs from the original
FP32-kernel objective, so the result is a useful low-precision execution
contract rather than an exact FP32 replacement.

The optimized pipeline reveals a more general systems lesson. Once regular
numeric work becomes inexpensive, the limiting resource can become dependent
optimization progress. The local pair-SMO solver does not lack arithmetic
throughput; each accepted step determines the next decision. Attempts to fill
the GPU by batching, replaying, or speculating can reduce visible depth while
increasing global-column generation, local-row consumption, or outer
iterations. Their success must therefore be measured by progress per weighted
kernel-data item, not occupancy alone.

Four principles emerge: preserve one coherent approximate operator across all
state-changing components; distinguish expensive global data production from
local data reuse; optimize the work--depth frontier rather than either axis in
isolation; and treat workload shape and solver phase as first-class dispatch
inputs. These principles are not specific to SVMs. They apply whenever an
iterative GPU algorithm alternates high-throughput data production with a
short, state-dependent decision chain. For RBF SVMs, the remaining opportunity
is a solver that discovers several jointly productive directions from each
working set, thereby obtaining substantially more convergence progress from
every generated kernel column without sacrificing the declared numerical
contract.
